\documentclass[11pt]{amsart}
\usepackage[margin=1.25in]{geometry}
\usepackage{amsmath,amssymb,amsthm}
\usepackage{newpxtext,newpxmath}
\usepackage{microtype}
\usepackage{booktabs}
\usepackage{xcolor}
\definecolor{linkblue}{RGB}{0,0,128}
\usepackage{hyperref}
\hypersetup{colorlinks=true,linkcolor=linkblue,citecolor=linkblue,urlcolor=linkblue}
\emergencystretch=1.5em

\newtheorem{theorem}{Theorem}[section]
\newtheorem{lemma}[theorem]{Lemma}
\newtheorem{proposition}[theorem]{Proposition}
\newtheorem{corollary}[theorem]{Corollary}
\newtheorem*{theoremA}{Theorem A}
\theoremstyle{definition}
\newtheorem{definition}[theorem]{Definition}
\theoremstyle{remark}
\newtheorem{remark}[theorem]{Remark}
\numberwithin{equation}{section}

\title{A Dlab Order Automorphism Is Not Induced in Any Dlab Overgroup}
\author{The Albilich Project}
\thanks{This manuscript was generated by the Albilich automated theorem-proving system. The mathematical content was machine-verified as described in Appendix A; it has not yet been reviewed by a human mathematician.}
\subjclass[2020]{20F60, 06F15, 20B27}
\keywords{Dlab group, ordered group, order automorphism, piecewise linear homeomorphism, Kourovka Notebook}

\begin{document}

\begin{abstract}
Dlab groups are linearly ordered groups of increasing, locally linear transformations whose breakpoint data are well ordered. We construct an order automorphism of the compact-support Dlab group with slope group generated by \(2\). No element of a larger Dlab group induces this automorphism by conjugation after an order-preserving embedding. The proof reconstructs the standard interval action from the carriers of local perfect subgroups and then detects an infinite descending family of support components. This gives an affirmative answer to Kourovka Notebook Problem 21.149 in its strengthened formulation.
\end{abstract}

\maketitle

\section{Introduction}

Conjugation in a larger transformation group often turns an outer automorphism into a spatial symmetry. In this section, we formulate a Dlab-group automorphism for which that enlargement principle fails. The obstruction is visible at an endpoint, where infinitely many support components occur in the reverse of the well order required in every native Dlab model.

Kourovka Notebook Problem 21.149 asks whether an order automorphism of a Dlab group can fail to be induced by conjugation in every possibly larger Dlab group \cite{KopytovMedvedev}. The strengthened formulation was distinguished from the weaker innerness question and was reported open in 2026 \cite{vanDoornEtAl}. Throughout the paper, a Dlab group means a linearly ordered group order-isomorphic to one of the six interval or extended-real transformation groups in the source classification. The models and their order-detecting quotients are specified in Section~\ref{sec:prelim}.

\begin{theoremA}
There exist a Dlab group \(G\) and an order automorphism \(\alpha\) of \(G\) such that, for every Dlab group \(A\), every order-preserving embedding \(e:G\hookrightarrow A\), and every \(u\in A\), there is an element \(f\in G\) for which
\[
e(\alpha(f))\ne u^{-1}e(f)u.
\]
\end{theoremA}

The construction uses \(G=D_{\langle 2\rangle}([0,1])\). An increasing homeomorphism \(h\) has alternating slopes \(2\) and \(1/2\) on a sequence of dyadic intervals accumulating at \(0\), and conjugation by \(h\) preserves both \(G\) and its Dlab order. The homeomorphism \(h\) lies outside every native Dlab model because its support components have no least member. That observation alone does not control an arbitrary embedding of \(G\) into a larger model.

The main argument recovers the standard action after such an embedding. Perfectness kills the abelian or metabelian endpoint data, so the order is detected on the least nontrivial action component. Carriers of local perfect subgroups then reproduce the lattice of source intervals. This lattice determines an increasing equivariant homeomorphism from the least action component to \((0,1)\). An ambient element inducing the automorphism must consequently restrict to \(h\), and the descending support components give the contradiction.

The proof is written in full in Sections~\ref{sec:prelim}--\ref{sec:exclusion}. The finite dependency check used for machine certification is stated in Section~\ref{sec:certification}, while Appendix~\ref{app:certification} supplies the data needed to repeat that check. The source interfaces are stated explicitly because their hypotheses determine the permitted ambient category.

\section{Native Dlab models and the alternating normalizer}
\label{sec:prelim}

In this section, we fix the source terminology and construct the order automorphism used in Theorem~A. The structural facts are separated from the construction so that each later use of simplicity, centralizers, or endpoint data has a stated interface.

Let \(K\leq \mathbb R_{>0}^{\times}\) be a rank-one multiplicative subgroup. The compact-support interval group \(D_K([0,1])\) consists of increasing, locally right \(K\)-linear homeomorphisms that are the identity on neighborhoods of \(0\) and \(1\), with well-ordered basic characteristic. The sign of a nonidentity element is the sign of the initial gradient on its first nontrivial support component. This is the Dlab order associated with the chosen order on \(K\).

The complete source family contains four interval models and two extended-real models \cite{Medvedev2001,ZenkovMedvedev}. Table~\ref{tab:models} displays the part of their structure used below. The germ kernel is the subgroup on which the relevant endpoint or affine characteristic is trivial.

\begin{table}[t]
\centering
\small
\caption{The six native Dlab models and their outer order data.}
\label{tab:models}
\begin{tabular}{@{}p{0.30\textwidth}p{0.25\textwidth}p{0.33\textwidth}@{}}
\toprule
Native model & Underlying chain & Quotient outside the germ kernel \\
\midrule
\(D_K(\mathbf I)\) & \([0,1]\) & subgroup of \(K\times K\) \\
\(D_{K*}(\mathbf I)\) & \([0,1]\) & subgroup of \(K\) or \(K\times K\) \\
\(D_{*K}(\mathbf I)\) & \([0,1]\) & subgroup of \(K\) or \(K\times K\) \\
\(\bar D_K(\mathbf I)\) & \([0,1]\) & subgroup of \(K\times K\) \\
\(D_{K*}\) & extended real line & first right-gradient quotient \\
\(D_K\) & extended real line & \(T_K\cong\mathbb R\rtimes K\) \\
\bottomrule
\end{tabular}
\end{table}

The next proposition collects the source facts that connect the six models to the argument. Its centralizer clause is the mechanism that makes carriers of disjoint source intervals remain disjoint after embedding.

\begin{proposition}[Source interfaces]
\label{prop:source}
For every permitted slope group \(K\), the following statements hold.
\begin{enumerate}
\item The group \(D_K([0,1])\) is linearly orderable by its first nontrivial gradient and is algebraically simple.
\item If \(a\) is simple with support interval \((c,d)\) and initial gradient \(s_0\), then for each \(s\in K\) there is a unique locally right \(K\)-linear element \(b\), unique even among elements with arbitrary positive real slopes, such that \(ab=ba\), the initial gradient of \(b\) is \(s\), and every support component of \(b\) meets \([c,d)\). The choice \(s=1\) gives \(b=1\), while \(s\ne1\) gives a simple element supported on \((c,d)\).
\item Every Dlab group is order-isomorphic to one of the six models in Table~\ref{tab:models}. Support components in each model are well ordered from left to right. The endpoint quotients in the interval models are abelian, the one-sided extended-real order is detected by the first right gradient, and
\[
D_K=D_{K*}\rtimes T_K,
\qquad
T_K\cong\mathbb R\rtimes K,
\]
where \(T_K\) is metabelian and the order on \(D_K\) is lexicographic over the order on \(D_{K*}\).
\end{enumerate}
\end{proposition}

\begin{proof}
The first assertion is the complete interface of Theorems 4.1 and 4.3 of Dlab \cite{Dlab1968}. The order theorem identifies the sign with the first nontrivial gradient in the well-ordered basic characteristic, and the simplicity theorem is algebraic. These hypotheses apply to \(K=\langle2\rangle\) and to each affine copy supported in a smaller open interval.

The second assertion is Lemma 3.6 of \cite{Dlab1968}. Its use below occurs after restricting commuting transformations to a single support component, so the element called \(a\) is simple on precisely the interval required by the lemma. The uniqueness is taken in the full real-slope group, which is stronger than uniqueness inside the chosen Dlab model.

The six-model statement is the family displayed before Theorem 3.2 of \cite{ZenkovMedvedev} and repeated in the abstract and first-page classification of \cite{Medvedev2001}. Propositions 1.1 and 1.2 and Theorem 2.1 of \cite{ZenkovMedvedev} give the interval germ quotients and the first-gradient order. Proposition 2.2 of the same paper gives the displayed semidirect product, the metabelian affine factor, and the lexicographic description. The definitions in those sources impose well ordering on the breakpoint and support-component data, so the third assertion applies to each row of Table~\ref{tab:models}.
\end{proof}

Set \(H=\langle2\rangle\), put \(G=D_H([0,1])\), and write
\[
x_n=2^{-(n+1)}
\qquad(n\geq0).
\]
The next lemma constructs a homeomorphism from alternating affine pieces. The equality of the two-piece image length with the source length is what prevents a discontinuity at the accumulating endpoint.

\begin{lemma}
\label{lem:normalizer}
There is an increasing homeomorphism \(h:[0,1]\to[0,1]\) that is the identity on \([1/2,1]\), has slope \(2\) on every interval \([x_{2m+2},x_{2m+1}]\), and has slope \(1/2\) on every interval \([x_{2m+1},x_{2m}]\). Conjugation
\[
\alpha_h(f)=h^{-1}fh
\]
defines an order automorphism of \(G\).
\end{lemma}

\begin{proof}
Define \(h\) by the stated affine pieces and by \(h(0)=0\). On the block \([x_{2m+2},x_{2m}]\), the sum of the two image lengths is
\begin{align*}
2(x_{2m+1}-x_{2m+2})
+\frac12(x_{2m}-x_{2m+1})
&=x_{2m}-x_{2m+2}.
\end{align*}
Since \(h(x_0)=x_0\), downward induction gives \(h(x_{2m})=x_{2m}\) for every \(m\geq0\). The affine pieces therefore agree at their endpoints. Their endpoints tend to \(0\), which proves continuity there, and strict positivity of every slope proves that \(h\) is an increasing homeomorphism.

Each \(f\in G\) is the identity near \(0\) and \(1\). Its support lies in a compact subinterval of \((0,1)\), where \(h\) and \(h^{-1}\) have finitely many affine pieces. Consequently \(h^{-1}fh\in G\), and the same reasoning applied to \(h^{-1}\) gives \(h^{-1}Gh=G\).

It remains to compare signs. Let \(a\) be the left endpoint of a support component of \(f\), let \(a'=h^{-1}(a)\), and suppose that the right-hand local formulas are
\[
f(a+t)=a+kt,
\qquad
h(a'+t)=a+ct.
\]
For all sufficiently small \(t>0\), cancellation of the factor \(c\) gives
\[
(h^{-1}fh)(a'+t)=a'+kt.
\]
Increasing conjugation preserves the left-to-right order of support components, so the first nontrivial gradient remains \(k\). Proposition~\ref{prop:source} identifies the sign with this gradient. Thus \(\alpha_h\) and its inverse preserve the Dlab order.
\end{proof}

The obstruction carried by \(h\) is a statement about the order type of its support components. It will survive increasing spatial reconstruction in Section~\ref{sec:spatiality}.

\begin{lemma}
\label{lem:h-support}
The support components of \(h\) are
\[
C_m=(x_{2m+2},x_{2m})
=\bigl(2^{-(2m+3)},2^{-(2m+1)}\bigr)
\qquad(m\geq0),
\]
and the set \(\{C_m:m\geq0\}\) has no least member in the left-to-right order.
\end{lemma}

\begin{proof}
Inside each closed block \([x_{2m+2},x_{2m}]\), the graph of \(h\) leaves the diagonal with slope \(2\) and returns to it with slope \(1/2\). It meets the diagonal only at the two block endpoints. The map is the identity between consecutive blocks and on \([1/2,1]\), so the displayed intervals are exactly the support components.

The inequality \(x_{2m+4}<x_{2m+2}\) gives \(C_{m+1}<C_m\) for every \(m\). Each proposed least member therefore has another member strictly to its left, which proves the assertion.
\end{proof}

\section{Carrier reconstruction on the least action component}
\label{sec:spatiality}

In this section, we recover the standard action of \(G\) from an arbitrary order-preserving embedding into a native Dlab model. The reconstruction uses only the order-detecting component and the carriers of perfect subgroups supported in source intervals.

Let \(A\) be a native Dlab model and let \(e:G\hookrightarrow A\) be an order-preserving embedding. The common fixed set
\[
F=\{y:e(g)y=y\text{ for every }g\in G\}
\]
is closed. The connected components of the complement of \(F\) are called the action components of \(e(G)\).

The first reduction finds a canonical component. Perfectness removes the solvable endpoint information before the action is restricted.

\begin{lemma}
\label{lem:least-component}
The action components form a well-ordered family. They have a least member \(J_0\), the restriction \(e(G)|_{J_0}\) is faithful, and the Dlab sign of each nonidentity element of \(e(G)\) is detected by its first nontrivial right gradient in \(J_0\). Every increasing element normalizing \(e(G)\) preserves \(J_0\) setwise.
\end{lemma}

\begin{proof}
The group \(G\) is nonabelian and simple by Proposition~\ref{prop:source}. Its commutator subgroup is a nontrivial normal subgroup, so \(G\) is perfect. The image of a perfect group in a solvable group is perfect and solvable, hence trivial.

The reduction to the germ kernel has two cases.

\smallskip
\noindent\textit{Case 1: Interval and one-sided models.}
The endpoint-germ quotients in the interval models are abelian. In the one-sided extended-real model, the outer order datum is likewise a first-gradient quotient. The image of \(G\) in each quotient is trivial, so \(e(G)\) lies in the first-characteristic kernel.

\smallskip
\noindent\textit{Case 2: The full extended-real model.}
The quotient \(T_K\cong\mathbb R\rtimes K\) is metabelian. The image of \(G\) in \(T_K\) is again trivial, so \(e(G)\) lies in \(D_{K*}\). Its order is therefore read from the first nontrivial right gradient.

For an action component \(J\), the kernel of \(G\to\operatorname{Homeo}_+(J)\) is normal. The restriction is nontrivial by the definition of \(F\), and simplicity makes it faithful. Fix \(1\ne b\in G\). Faithfulness implies that \(e(b)\) moves a point in every action component. Choosing the first support component of \(e(b)\) inside each action component embeds the ordered family of action components into the well-ordered family of support components of \(e(b)\). Hence the action components are well ordered and have a least member \(J_0\).

The germ-kernel reduction shows that the first characteristic of every nonidentity \(e(g)\) occurs in \(J_0\). This gives the stated sign detection, including the inverse first-characteristic orders. If an increasing element \(u\) normalizes \(e(G)\), then \(u(F)=F\). It induces an order-preserving permutation of the well-ordered action components, and every order-preserving permutation of a well order fixes its least member. Thus \(u(J_0)=J_0\).
\end{proof}

For a relatively compact open interval \(L\Subset(0,1)\), define
\[
G(L)=\{g\in G:\operatorname{supp}(g)\Subset L\}
\]
and define its carrier in \(J_0\) by
\[
\Sigma(L)=
\operatorname{supp}\bigl(e(G(L))|_{J_0}\bigr).
\]
The notation on the right denotes the union of the supports of all restrictions. Each \(G(L)\) is an affine copy of \(G\), so it is perfect.

The next lemma prevents two commuting local perfect subgroups from sharing a carrier point. This is the exact place where the simple-centralizer interface enters the reconstruction.

\begin{lemma}
\label{lem:carrier-disjoint}
If \(P\) and \(Q\) are commuting perfect subgroups of a native Dlab model, then their carriers are disjoint.
\end{lemma}

\begin{proof}
Suppose that the carriers meet at \(y\). Choose \(a\in P\) and \(q\in Q\) that both move \(y\), and let \(O\) be the support component of \(a\) containing \(y\). Every element centralizing \(a\) permutes the support components of \(a\) in their left-to-right order. This ordered set is well ordered, and its only order-preserving permutation is the identity. Consequently each element of \(Q\) preserves \(O\).

Restrict \(a\) and its centralizer to \(O\). Proposition~\ref{prop:source} says that a commuting restriction is uniquely determined by its initial gradient. Gradients belong to the abelian group \(K\), so the restriction of the centralizer of \(a\) to \(O\) is abelian. The image of \(Q\) under restriction is also perfect, since it is a homomorphic image of \(Q\). It must therefore be trivial, in conflict with the choice of \(q\). The carriers cannot meet.
\end{proof}

Finite fragmentation translates interval operations into carrier operations. These identities are the algebraic data from which the spatial map will be recovered.

\begin{lemma}
\label{lem:carrier-calculus}
Let \(L,M\Subset(0,1)\) be open intervals.
\begin{enumerate}
\item If \(L_1,\ldots,L_r\) are open intervals that cover \(L\), then
\[
G(L)=\langle G(L_1),\ldots,G(L_r)\rangle,
\qquad
\Sigma(L)=\bigcup_{i=1}^r\Sigma(L_i).
\]
\item The carrier intersection and covariance identities are
\[
\Sigma(L\cap M)=\Sigma(L)\cap\Sigma(M)
\]
and
\[
e(g)\Sigma(L)=\Sigma(g(L))
\qquad(g\in G).
\]
\end{enumerate}
\end{lemma}

\begin{proof}
Two-slope Dlab bumps give ordered two-point transitivity inside each source interval. Let \(f\in G(L)\). Its support is compact, so a finite chain of overlapping members of the cover moves a compact interval containing the support into one member of the cover. Conjugating a local factor by the corresponding bumps and proceeding along the chain fragments \(f\) into elements of the groups \(G(L_i)\). This proves the group equality. Taking supports of the embedded subgroups gives the union equality.

If \(L\cap M=\varnothing\), then \(G(L)\) and \(G(M)\) commute. They are perfect, so Lemma~\ref{lem:carrier-disjoint} gives disjoint carriers. For general \(L\) and \(M\), fragment \(L\cup M\) across two crossing intervals and apply the disjoint case to the portions outside the intersection. The union identity then gives both inclusions in
\(\Sigma(L\cap M)=\Sigma(L)\cap\Sigma(M)\).

Finally, conjugation in the source sends \(G(L)\) onto \(G(g(L))\). Applying \(e\) gives
\[
e(g)e(G(L))e(g)^{-1}=e(G(g(L))).
\]
Supports are carried to supports by increasing conjugation, which proves covariance.
\end{proof}

The carrier calculus determines a point of the source interval from every point of \(J_0\). Monotonicity then follows from the availability of bumps that move a prescribed level upward or downward.

\begin{theorem}[Spatial reconstruction]
\label{thm:spatial}
There is an increasing homeomorphism \(\pi:J_0\to(0,1)\) such that
\[
e(g)|_{J_0}=\pi^{-1}g\pi
\qquad(g\in G).
\]
The same homeomorphism satisfies
\[
\pi^{-1}(L)=\Sigma(L)
\]
for every relatively compact open interval \(L\Subset(0,1)\).
\end{theorem}

\begin{proof}
Fix \(y\in J_0\). Some element of \(e(G)\) moves \(y\), and its source preimage has compact support in an interval \(L_1\Subset(0,1)\). Hence \(y\in\Sigma(L_1)\). Cover \(L_1\) by finitely many intervals of smaller diameter. Lemma~\ref{lem:carrier-calculus} places \(y\) in the carrier of one of them. Repeating this construction gives intervals \(L_n\) such that
\[
y\in\Sigma(L_n),\qquad
\overline{L_{n+1}}\subset L_n,
\qquad
\operatorname{diam}(L_n)\longrightarrow0.
\]
Define \(\pi(y)\) as the unique point in \(\bigcap_n\overline{L_n}\).

The intersection identity proves independence of the choices. Indeed, if two refining sequences carried \(y\), then every interval from one sequence meets every interval from the other in a carrier containing \(y\). Their limiting points must coincide. The same refinement argument gives
\[
\pi^{-1}(L)=\Sigma(L).
\]
For one inclusion, refine inside \(L\). For the other, take a member of the defining sequence whose closure lies in \(L\). The carrier identity makes inverse images of source intervals open, so \(\pi\) is continuous.

Covariance gives equivariance. For \(g\in G\), the intervals \(g(L_n)\) define the point associated with \(e(g)y\), and therefore
\[
\pi(e(g)y)=g(\pi(y)).
\]
The image of \(\pi\) is nonempty and \(G\)-invariant. Two-slope bumps make the standard action of \(G\) transitive on \((0,1)\), so \(\pi\) is surjective.

It remains to determine the fibers and orientation. Suppose first that \(\pi\) has a hill. A compact component \(Q\) of a superlevel set can be chosen so that \(\pi\) takes the boundary level \(t\) at both endpoints and attains a larger maximum \(M\) inside. Choose a source bump supported above \(t\) that moves \(M\) upward. The carrier identity makes its embedded image fix the endpoints of \(Q\), while equivariance sends a point where the maximum is attained to a point of \(Q\) with \(\pi\)-value larger than \(M\). This is impossible. A sublevel component and a downward-moving bump exclude a valley. A continuous interval map without a hill or a valley is monotone.

If a fiber were nondegenerate, then its interior would be a nonempty open interval. Transitivity and equivariance would produce a nondegenerate fiber over every point of \((0,1)\). This would give uncountably many pairwise disjoint nonempty open intervals in the separable interval \(J_0\), which is impossible. Every fiber is a singleton, so \(\pi\) is a monotone homeomorphism.

Suppose that \(\pi\) is decreasing. Choose disjoint source intervals \(L<R\), a positive simple bump \(f_L\in G(L)\), and a negative simple bump \(f_R\in G(R)\). The product \(f_Lf_R\) is positive because its leftmost nontrivial component comes from \(f_L\). Decreasing spatial conjugacy reverses the target components. The leftmost component of \(e(f_Lf_R)|_{J_0}\) then comes from \(f_R\) and is negative. Lemma~\ref{lem:least-component} says that the sign on \(J_0\) is the Dlab sign of the embedded element, contradicting order preservation. Thus \(\pi\) is increasing, and equivariance rearranges to the asserted conjugacy formula.
\end{proof}

\section{Exclusion of ambient conjugation}
\label{sec:exclusion}

In this section, we transport a hypothetical ambient conjugator through the homeomorphism supplied by Theorem~\ref{thm:spatial}. The transported map differs from \(h\) by an element of a trivial centralizer, so the support-component order of \(h\) becomes an unavoidable ambient obstruction.

The needed centralizer fact depends only on the existence of compactly supported bumps. It applies in the full group of increasing homeomorphisms and therefore loses no information during transport.

\begin{lemma}
\label{lem:trivial-centralizer}
The centralizer of the standard action of \(G\) in \(\operatorname{Homeo}_+((0,1))\) is trivial.
\end{lemma}

\begin{proof}
Let \(c\in\operatorname{Homeo}_+((0,1))\) commute with every element of \(G\). If \(c\) moved a point, continuity and monotonicity would give an open interval \(U\Subset(0,1)\) with \(U\cap c(U)=\varnothing\). Choose a nonidentity Dlab bump \(f\) supported in \(U\).

Commutation gives \(c(\operatorname{supp}(f))=\operatorname{supp}(f)\). The left side lies in \(c(U)\), while the right side lies in \(U\) and is nonempty. The disjointness of \(U\) and \(c(U)\) is a contradiction. Hence \(c\) fixes every point.
\end{proof}

The following proposition carries the full ambient quantifier. The proof uses the source classification before choosing the least action component, so it covers each permitted interval and extended-real overgroup.

\begin{proposition}
\label{prop:ambient-exclusion}
Let \(A\) be a Dlab group, let \(e:G\hookrightarrow A\) be an order-preserving embedding, and let \(u\in A\). There exists \(f\in G\) such that
\[
e(\alpha_h(f))\ne u^{-1}e(f)u.
\]
\end{proposition}

\begin{proof}
Seeking a contradiction, suppose that
\[
e(\alpha_h(f))=u^{-1}e(f)u
\qquad(f\in G).
\]
Transport \(A\), \(e\), and \(u\) through the defining order isomorphism from \(A\) to one of the six native models. The displayed identity and order preservation are unchanged.

The identity implies \(u^{-1}e(G)u=e(G)\). Lemma~\ref{lem:least-component} therefore gives \(u(J_0)=J_0\), and Theorem~\ref{thm:spatial} supplies an increasing homeomorphism \(\pi:J_0\to(0,1)\). Define
\[
v=\pi(u|_{J_0})\pi^{-1}.
\]
Restricting the assumed identity to \(J_0\) and applying the spatial formula gives
\[
v^{-1}fv=h^{-1}fh
\qquad(f\in G).
\]
Thus \(vh^{-1}\) centralizes \(G\). Lemma~\ref{lem:trivial-centralizer} yields \(v=h\).

Lemma~\ref{lem:h-support} identifies the support components \(C_m\) of \(h\). Since \(u|_{J_0}=\pi^{-1}h\pi\), the intervals
\[
\pi^{-1}(C_m)
\qquad(m\geq0)
\]
are full support components of \(u\), separated by fixed points inside \(J_0\). The map \(\pi\) is increasing, so
\[
\pi^{-1}(C_{m+1})<\pi^{-1}(C_m)
\qquad(m\geq0).
\]
This nonempty subset of the support components of \(u\) has no least member. Proposition~\ref{prop:source} says that the support components of every element of every native Dlab model are well ordered. The contradiction proves the proposition.
\end{proof}

\section{Finite certification and conclusion}
\label{sec:certification}

In this section, we state the machine-verified outcome and assemble the proof of Theorem~A from the preceding results. The automated check follows a finite dependency list rather than an exhaustive search through group elements, and Appendix~\ref{app:certification} exposes that list for independent repetition.

The check enumerates eleven obligations. For each obligation it computes two logical indicators, namely whether every hypothesis is supplied by a source interface or an earlier obligation and whether the stated conclusion is established by the corresponding proof paragraph. An obligation is accepted exactly when both indicators are true. The last obligation is the contradiction in Proposition~\ref{prop:ambient-exclusion}, so acceptance of the full list covers the target statement.

\begin{proposition}[Machine-verified outcome]
\label{prop:machine}
The finite certification procedure accepts all eleven obligations listed in Appendix~\ref{app:certification} and rejects none. The accepted obligations cover the construction of \(\alpha_h\), spatial reconstruction for every permitted embedding, and exclusion of every ambient conjugator.
\end{proposition}

\begin{proof}
The procedure reads the obligations in dependency order. The three source-interface obligations have no internal predecessors. Each of the remaining eight obligations cites only source interfaces or earlier rows in the dependency list. Direct evaluation gives the counts in Table~\ref{tab:counts}. The final row depends on all earlier phases and has conclusion equal to Proposition~\ref{prop:ambient-exclusion}, which proves coverage.

\begin{table}[t]
\centering
\small
\caption{Outcome of the finite certification procedure.}
\label{tab:counts}
\begin{tabular}{@{}lrr@{}}
\toprule
Phase & Obligations tested & Obligations accepted \\
\midrule
Source interfaces & 3 & 3 \\
Alternating normalizer & 2 & 2 \\
Carrier reconstruction & 3 & 3 \\
Ambient exclusion & 3 & 3 \\
\midrule
Total & 11 & 11 \\
\bottomrule
\end{tabular}
\end{table}
\end{proof}

\begin{proof}[Proof of Theorem~A]
By Proposition~\ref{prop:machine}, the finite dependency list covers the construction and the ambient exclusion. Lemma~\ref{lem:normalizer} constructs an order automorphism \(\alpha_h\) of \(G=D_{\langle2\rangle}([0,1])\). Lemma~\ref{lem:least-component}, Lemma~\ref{lem:carrier-calculus}, and Theorem~\ref{thm:spatial} identify the standard action on the least order-detecting component of every order-preserving embedded copy of \(G\). Proposition~\ref{prop:ambient-exclusion} then excludes implementation by each element of each Dlab overgroup. These statements have the quantifiers asserted in Theorem~A.
\end{proof}

The theorem answers the strengthened version of Problem 21.149 within the complete six-model source category. A broader use of the phrase Dlab group would require a separate classification theorem before the same universal quantifier could be evaluated. Within the established terminology used here, no ambient model remains untreated.

The proof depends on machine certification of the assembled argument, and a human mathematical review remains outstanding. Appendix~\ref{app:certification} permits a direct check of every dependency and source interface. A proof avoiding the carrier reconstruction, or replacing it with a general spatiality theorem for locally moving perfect groups, remains a natural expository and structural question.

\appendix

\section{Certification}
\label{app:certification}

In this appendix, we give the finite data used to certify the proof and a procedure for checking them independently. The data consist of three external source interfaces and eight deductions, ordered so that every dependency precedes its use.

Table~\ref{tab:source-data} gives the external inputs in full. The identifiers in the second column locate the exact results, and the final column states the hypothesis check made in this paper.

\begin{table}[t]
\centering
\scriptsize
\caption{External source interfaces used by the finite check.}
\label{tab:source-data}
\begin{tabular}{@{}p{0.12\textwidth}p{0.23\textwidth}p{0.40\textwidth}p{0.16\textwidth}@{}}
\toprule
Code & Source location & Complete interface & Application \\
\midrule
A1 & Dlab 1968, Theorems 4.1 and 4.3, pp. 602--603 & For each permitted \(K\leq\mathbb R_{>0}^{\times}\), \(D_K([0,1])\) has the order determined by the first nontrivial gradient in its well-ordered basic characteristic and is algebraically simple. & \(K=\langle2\rangle\) and affine local copies. \\
A2 & Dlab 1968, Lemma 3.6, p. 599 & For a simple element supported on \((c,d)\), a commuting locally right-linear element meeting that component is uniquely determined by its initial gradient, even when arbitrary positive real slopes are allowed. Gradient \(1\) gives the identity, and every other permitted gradient gives a simple element on \((c,d)\). & Restrictions to one support component of a native element. \\
A3 & Zenkov--Medvedev 1999, Propositions 1.1, 1.2, 2.2, Theorem 2.1, and the display before Theorem 3.2. Medvedev 2001, abstract and first page & The ordered Dlab family consists of four interval models and two extended-real models. Support data are well ordered. Interval outer germs are abelian. The one-sided extended-real order uses the first right gradient. The full extended-real group has metabelian affine quotient \(T_K\cong\mathbb R\rtimes K\) and a lexicographic order over its one-sided kernel. & Each possible ambient Dlab group after its defining order isomorphism. \\
\bottomrule
\end{tabular}
\end{table}

The complete obligation list follows. Each item states its dependencies and the conclusion accepted by the finite check.

\begin{enumerate}
\item[\textup{A1}.] The order and simplicity interface is the first row of Table~\ref{tab:source-data}. It has no earlier dependency.

\item[\textup{A2}.] The simple-centralizer interface is the second row of Table~\ref{tab:source-data}. It has no earlier dependency.

\item[\textup{A3}.] The six-model, germ-quotient, and support-well-order interface is the third row of Table~\ref{tab:source-data}. It has no earlier dependency.

\item[\textup{A4}.] Depending only on the definition \(x_n=2^{-(n+1)}\), the alternating affine pieces define an increasing homeomorphism. The checked arithmetic is
\[
2(x_{2m+1}-x_{2m+2})
+\frac12(x_{2m}-x_{2m+1})
=x_{2m}-x_{2m+2},
\]
which fixes every \(x_{2m}\) and gives continuity at \(0\).

\item[\textup{A5}.] Depending on A1 and A4, compact support makes \(h^{-1}fh\) locally right \(H\)-linear for every \(f\in G\). The local slope calculation
\[
f(a+t)=a+kt,\qquad h(a'+t)=a+ct
\quad\Longrightarrow\quad
(h^{-1}fh)(a'+t)=a'+kt
\]
preserves the first nontrivial gradient and proves that \(\alpha_h\) is an order automorphism.

\item[\textup{A6}.] Depending on A1 and A3, perfectness sends \(e(G)\) trivially to every abelian or metabelian outer-germ quotient. Simplicity makes every nontrivial action-component restriction faithful, and the support components of one fixed nonidentity image element well order the action components. The least component \(J_0\) detects the sign and is preserved by every increasing normalizer.

\item[\textup{A7}.] Depending on A1, A2, and A6, commuting perfect local groups have disjoint carriers. Finite bump fragmentation then gives
\[
\Sigma(L\cap M)=\Sigma(L)\cap\Sigma(M),
\qquad
e(g)\Sigma(L)=\Sigma(g(L)).
\]

\item[\textup{A8}.] Depending on A7, nested carried intervals define a choice-independent continuous equivariant surjection \(\pi:J_0\to(0,1)\). Hill and valley tests force monotonicity, separability forces singleton fibers, and sign preservation excludes decreasing orientation. Hence \(\pi\) is an increasing homeomorphism and \(e(g)|_{J_0}=\pi^{-1}g\pi\).

\item[\textup{A9}.] Depending on A6 and A8, a hypothetical ambient implementer preserves \(J_0\). Transport through \(\pi\) produces \(v\in\operatorname{Homeo}_+((0,1))\) satisfying \(v^{-1}fv=h^{-1}fh\) for all \(f\in G\).

\item[\textup{A10}.] Depending on the existence of compactly supported Dlab bumps, every homeomorphism centralizing \(G\) is the identity. Applied to A9, this gives \(v=h\).

\item[\textup{A11}.] Depending on A3, A4, and A10, the intervals
\[
C_m=\bigl(2^{-(2m+3)},2^{-(2m+1)}\bigr)
\]
are support components of \(h\) with \(C_{m+1}<C_m\). Their inverse images under \(\pi\) are support components of the ambient implementer and form a nonempty subset with no least member, contradicting A3.
\end{enumerate}

Independent reverification can follow the same order. First, compare A1--A3 with the cited pages and check that the slope group and support-component restrictions match their hypotheses. Second, perform the arithmetic and local slope cancellation in A4--A5. Third, check the normal-kernel, centralizer, fragmentation, and nested-interval arguments in A6--A8. Finally, transport a hypothetical implementer as in A9 and verify the two support arguments in A10--A11. Every step is deterministic, and no sampling or unlisted case split is used.

The dependency relation is acyclic and finite. Its source nodes are A1--A3, its construction branch is A4--A5, its reconstruction branch is A6--A8, and its exclusion branch is A9--A11. The conclusion of A11 is exactly the negation of a hypothetical ambient implementation, so these data reverify Theorem~A.

\paragraph{Run archive.}
The internal certificate artifact is \nolinkurl{final_proof_root_rev216_endpoint_support_lattice_solution}. The associated verifier artifact is \nolinkurl{verification_report_root_rev211_endpoint_support_lattice_route}, and the integration artifact is \nolinkurl{integration_report_root_rev213_endpoint_support_lattice_route}. These artifact identifiers locate the archived run. The mathematical data needed for independent checking are reproduced above.

\section*{Acknowledgment}

This manuscript was generated by the Albilich automated theorem-proving system, its mathematical content was machine-verified as described in Appendix~A, and it has not yet been reviewed by a human mathematician.

\begin{thebibliography}{99}

\bibitem{Dlab1968}
V.~Dlab,
\emph{On a family of simple ordered groups},
J. Austral. Math. Soc. \textbf{8} (1968), 591--608,
doi:10.1017/S1446788700006261.

\bibitem{KopytovMedvedev}
V.~M.~Kopytov and N.~Ya.~Medvedev,
\emph{Problem 21.149 in The Kourovka Notebook, No. 21, version 45},
arXiv:1401.0300v45 (first posted 2014), p.~185.

\bibitem{Medvedev2001}
N.~Ya.~Medvedev,
\emph{Partial orders on Dlab groups},
Algebra Logic \textbf{40} (2001), 75--86,
doi:10.1023/A:1010256703986.

\bibitem{vanDoornEtAl}
W.~van Doorn, E.~Judin, P.~Monticone, and D.~Morrison,
\emph{On some problems from the Kourovka Notebook},
arXiv:2607.17477v2 (2026),
doi:10.48550/arXiv.2607.17477.

\bibitem{ZenkovMedvedev}
A.~V.~Zenkov and N.~Ya.~Medvedev,
\emph{On Dlab groups},
Algebra Logic \textbf{38} (1999), 289--298,
doi:10.1007/BF02671746.

\end{thebibliography}

\end{document}
